% Generated draft: requires mathematical and source review.
\documentclass[11pt]{article}
\usepackage[margin=25mm]{geometry}
\usepackage{fontspec}
\setmainfont{Latin Modern Roman}
\usepackage{amsmath,amssymb,mathtools,unicode-math}
\setmathfont{Latin Modern Math}
\usepackage{enumitem,needspace,xurl,hyperref}
\setlength{\parindent}{0pt}
\setlength{\parskip}{5pt}
\setlength{\emergencystretch}{3em}
\newcommand{\N}{\mathbb{N}}
\newcommand{\Z}{\mathbb{Z}}
\newcommand{\Q}{\mathbb{Q}}
\newcommand{\R}{\mathbb{R}}
\newcommand{\C}{\mathbb{C}}
\newcommand{\F}{\mathbb{F}}
\newcommand{\A}{\mathbb{A}}
\newcommand{\PP}{\mathbb{P}}
\newcommand{\E}{\mathbb{E}}
\newcommand{\eps}{\varepsilon}
\newcommand{\dd}{\,\mathrm{d}}
\begin{document}
\section*{UnsolvedMath ID: 701669}
\textbf{Source problem code:} catalogue entry
\subsection{Definitions and mathematical statement}
Let \(F\leq H\) be free groups, with no assumption that either has finite rank. For a free group \(K\), let \(\operatorname{rk}(K)\) denote the cardinality of a free basis; the trivial group has rank \(0\). Assume that
\[
\mathcal G(F,H)=\{G\mid F\leq G\leq H,\ \operatorname{rk}(G)<\infty\}
\]
is nonempty. Define
\[
r=r(F,H)=\min\{\operatorname{rk}(G):G\in\mathcal G(F,H)\}\in\Z_{\geq0},
\qquad
\mathcal M(F,H)=\{G\in\mathcal G(F,H):\operatorname{rk}(G)=r\}.
\]
For subgroups \(G_1,G_2\leq H\), write \(\langle G_1,G_2\rangle\) for the subgroup of \(H\) generated by their union.

Determine which, if any, of the following assertions hold for every pair \(F\leq H\) satisfying these hypotheses:
\begin{enumerate}
\item[(i)] There exists \(G_{\max}\in\mathcal M(F,H)\) such that \(G\leq G_{\max}\) for every \(G\in\mathcal M(F,H)\).
\item[(i\(^{\prime}\))] For every \(G_1,G_2\in\mathcal M(F,H)\),
\[
\operatorname{rk}(\langle G_1,G_2\rangle)=r.
\]
\item[(ii)] There exists \(G_{\min}\in\mathcal M(F,H)\) such that \(G_{\min}\leq G\) for every \(G\in\mathcal M(F,H)\).
\item[(ii\(^{\prime}\))] For every \(G_1,G_2\in\mathcal M(F,H)\),
\[
\operatorname{rk}(G_1\cap G_2)=r.
\]
\end{enumerate}
Here largest and smallest mean greatest and least with respect to subgroup inclusion, respectively.

For each displayed label \(j\), let \(P_j\) denote the assertion that condition \(j\) holds for every admissible pair \(F\leq H\). The supplied source additionally asserts
\[
P_{\mathrm{(i)}}\Longrightarrow P_{\mathrm{(i^{\prime})}},\qquad
P_{\mathrm{(ii)}}\Longrightarrow P_{\mathrm{(ii^{\prime})}},\qquad
P_{\mathrm{(i^{\prime})}}\Longrightarrow P_{\mathrm{(i)}}.
\]
These are retained as source commentary, without proof or independent verification. For the last implication the source gives the explanation that ``subgroups of a free group of any fixed finite rank satisfy ACC,'' where ACC means the ascending chain condition. The wording does not explicitly specify whether the fixed rank concerns the subgroups in a chain or their ambient free group; this qualification remains unresolved in this draft. The source also reports that the proposer did not see how to obtain the converse \(P_{\mathrm{(ii^{\prime})}}\Longrightarrow P_{\mathrm{(ii)}}\); that report does not assert that this converse is false.
\subsection{Short English statement}
Among all finite-rank free groups lying between \(F\) and \(H\), consider those of the least possible rank \(r\). Must this collection have a greatest member or a least member under inclusion? Must the subgroup generated by any two members, or their intersection, also have rank \(r\)? Each question is quantified over every admissible inclusion \(F\leq H\). The supplied record labels the problem open; this status and its accompanying implication claims have not been independently verified.
\subsection{Sources}
\begin{itemize}
\item Local UnsolvedMath registry, permanent ID \(701669\), exact source problem code \texttt{catalogue entry}, title \emph{Kourovka Notebook Problem 20.11}. The supplied record was created and updated on \texttt{2026-10-16T14:00:00Z}; upstream revision \texttt{d371ee917ac9647fbb06cd6fc529deec6f3dda82}. Its supplied status is \texttt{open}; this is registry metadata, not an independently verified mathematical assessment.
\item The supplied attribution is to G. M. Bergman, proposed in 2022. The background identifies \emph{The Kourovka Notebook (Unsolved Problems in Group Theory)}, 20th issue (2022), original problem number 20.11, and the 21st-issue edition (2026), identified as \texttt{arXiv:1401.0300v48}.
\item The supplied collection description credits E. I. Khukhro and V. D. Mazurov as editors of the \emph{Kourovka Notebook}, Issues 1--21 (1965--2026). The bibliographic attribution and explanatory comments above are reproduced from the supplied record; the cited editions were not independently consulted.
\end{itemize}
\end{document}
